\begin{figure}[H]
\centering
\begin{tikzpicture}[
    font=\small,
    >=Latex,
    node distance=7mm and 15mm,
    block/.style={rectangle, rounded corners=2mm, draw, thick, align=center,
        text width=42mm, minimum height=9mm},
    runtime/.style={block, draw=gamegreen, fill=gamegreen!8},
    symbolic/.style={block, draw=htnblue, fill=htnblue!8},
    control/.style={block, draw=fallbackorange, fill=fallbackorange!10},
    decision/.style={diamond, aspect=2.1, draw=black!70, thick, align=center,
        fill=black!3, text width=29mm, inner sep=1.5pt},
    arrow/.style={->, thick}
]
    \node[runtime] (env) {Ambiente GridWorld};
    \node[runtime, below=of env] (sensor) {Sensor atualiza fatos simbólicos};
    \node[symbolic, below=of sensor] (state) {\texttt{WorldState} atualizado};
    \node[control, below=of state] (update) {\textit{Delegate} notifica o agente, que atualiza o planner};
    \node[control, below=of update] (validate) {Validar o plano remanescente no próximo tick};
    \node[decision, below=of validate] (valid) {O plano continua válido?};
    \node[symbolic, below left=of valid] (keep) {Manter execução do plano};
    \node[control, below right=of valid] (replan) {Solicitar novo plano HTN};

    \draw[arrow, gamegreen] (env) -- (sensor);
    \draw[arrow, htnblue] (sensor) -- (state);
    \draw[arrow, fallbackorange] (state) -- (update);
    \draw[arrow, fallbackorange] (update) -- (validate);
    \draw[arrow] (validate) -- (valid);
    \draw[arrow, htnblue] (valid) -- node[above left] {sim} (keep);
    \draw[arrow, fallbackorange] (valid) -- node[above right] {não} (replan);
\end{tikzpicture}
\caption{Atualização sensorial e replanejamento preguiçoso. Uma mudança de
estado não descarta imediatamente o plano: o agente valida o plano remanescente
no tick seguinte e só solicita novo planejamento se ele se tornou inválido.}
\label{fig:sensores-replanejamento}
\end{figure}
